\section*{Capítulo \ref{ch:b3:bacteriology} --- Bacteriologia: crescimento, fisiologia e genética}

\begin{solution}{exo:b3:bacteriology:1}
\omterm{def:b3:bacteriology:envelope}{Gram-positiva}: membrana citoplasmática envolvida por uma parede de
\omterm{def:b3:bacteriology:envelope}{peptidoglicano} espessa, de muitas camadas, entremeada de ácidos
teicoicos. \omterm{def:b3:bacteriology:envelope}{Gram-negativa}: uma camada fina de \omterm{def:b3:bacteriology:envelope}{peptidoglicano} num
periplasma e, depois, uma \omterm{def:b3:bacteriology:envelope}{membrana externa} cuja monocamada externa é
\omterm{def:b3:bacteriology:envelope}{lipopolissacarídeo}, com porinas. O complexo cristal violeta--iodo fica
preso na parede espessa e desidratada pelo álcool da \omterm{def:b3:bacteriology:envelope}{Gram-positiva}
(roxa), mas é lavado para fora da \omterm{def:b3:bacteriology:envelope}{Gram-negativa} por sua parede fina e
sua \omterm{def:b3:bacteriology:envelope}{membrana externa} rompida, que então absorve a contracoloração rosa.
\end{solution}

\begin{solution}{exo:b3:bacteriology:2}
$10^{6}$ vezes $= 2^{19.9}$: cerca de $20$ gerações;
$T = \qty{8}{h}/20
= \qty{24}{min}$;
$\mu = \ln 2/\qty{0.4}{h} = \qty{1.7}{h^{-1}}$.
\end{solution}

\begin{solution}{exo:b3:bacteriology:3}
\omterm{def:b3:bacteriology:hgt}{Transformação} --- Griffith (1928) e Avery, MacLeod e McCarty (1944):
pneumococos virulentos mortos pelo calor transformam pneumococos vivos
avirulentos, e o agente é o DNA. \omterm{def:b3:bacteriology:hgt}{Transdução} --- Zinder e Lederberg
(1952): o fago P22 leva genes de \emph{Salmonella} de uma linhagem a
outra através de um filtro que retém as células. \omterm{def:b3:bacteriology:hgt}{Conjugação} ---
Lederberg e Tatum (1946): duas linhagens auxotróficas de
\emph{E.~coli} dão recombinantes prototróficos apenas quando
misturadas, exigindo contato celular (o tubo em U de Davis).
\end{solution}

\begin{solution}{exo:b3:bacteriology:4}
Penicilina: as transpeptidases que fazem as ligações cruzadas do
\omterm{def:b3:bacteriology:envelope}{peptidoglicano}, que as células animais não têm. Estreptomicina: a
subunidade ribossômica pequena, cuja forma bacteriana difere da
eucarionte. Ciprofloxacina: a DNA-girase, uma topoisomerase bacteriana
ausente nos humanos. Rifampicina: a subunidade $\beta$ da
RNA-polimerase bacteriana, que a enzima eucarionte não compartilha.
\end{solution}

\begin{solution}{exo:b3:bacteriology:5}
$S^{*} = K_{S}D/(\mu_{\max} - D)$, $X^{*} = Y(S_{0} - S^{*})$.
$D = 0.2$: $S^{*} = \qty{0.0067}{g/L}$, $X^{*} = \qty{0.80}{g/L}$.
$D = 0.6$: $S^{*} = 0.06$, $X^{*} = 0.78$. $D = 0.79$: $S^{*} = 1.58$,
$X^{*} = 0.17$. Lavagem quando
$D = \mu(S_{0}) = 0.8\times 2/2.02 = \qty{0.79}{h^{-1}}$.
\end{solution}

\begin{solution}{exo:b3:bacteriology:6}
$(N/V)_{\text{thr}} = kA_{\text{thr}}/p$. Órgão de luz: $0.01\times
6\times 10^{12}/100 = 6\times 10^{8}$ por litro $= 6\times 10^{5}$ por
mL. Água do mar: $10\times 6\times 10^{12}/100 = 6\times 10^{11}$ por
litro $= 6\times 10^{8}$ por mL --- mil vezes mais alta, nunca alcançada
no mar.
\end{solution}

\begin{solution}{exo:b3:bacteriology:7}
\qty{25}{mm}: sensível, CIM baixa. \qty{12}{mm}: suscetibilidade
reduzida, intermediária --- pode falhar em doses comuns. \qty{0}{mm}:
resistente, com crescimento até o disco. O fármaco se difunde do disco
de modo que sua concentração cai com a distância de forma
reprodutível; o crescimento para onde a concentração iguala a CIM, de
modo que o raio da zona e a CIM se relacionam por uma curva de
calibração estabelecida com linhagens de CIM conhecida.
\end{solution}

\begin{solution}{exo:b3:bacteriology:8}
O \omterm{def:b3:genetic-engineering:tools}{plasmídeo} conjugativo carrega um integron cuja série de cassetes
contém vários genes de resistência, além de transpósons que carregam
mais; o \omterm{def:b3:genetic-engineering:tools}{plasmídeo} cruza num único acasalamento e expressa todos eles.
Eles se agrupam porque a seleção por qualquer um dos fármacos mantém o
\omterm{def:b3:genetic-engineering:tools}{plasmídeo} inteiro, de modo que os genes que pegam carona nele persistem
(cosseleção); os integrons capturam cassetes; os transpósons saltam
para \omterm{def:b3:genetic-engineering:tools}{plasmídeos}; e a unidade de transferência é o \omterm{def:b3:genetic-engineering:tools}{plasmídeo}, e não o
gene.
\end{solution}

\begin{solution}{exo:b3:bacteriology:9}
Só o fármaco: $10^{-7}\times 10^{9} = 100$ células resistentes
esperadas. Dois fármacos independentes: $10^{-14}\times 10^{9} = 10^{-5}$. As persistentes não são mutantes: uma pequena fração de
células dormentes sobrevive a qualquer fármaco e volta a crescer numa
população plenamente sensível quando o fármaco se vai, de modo que um
segundo fármaco não ajuda; só ajuda um tratamento longo o bastante para
apanhá-las quando despertam, ou um fármaco ativo em células dormentes.
\end{solution}

\begin{solution}{exo:b3:bacteriology:10}
$X > X^{*}$: o consumo $\mu X/Y$ excede o suprimento, e por isso $S$
cai abaixo de $S^{*}$; então $\mu(S) < D$ e
$\mathrm{d}X/\mathrm{d}t < 0$, e $X$ volta a cair; à medida que $X$
cai, $S$ sobe de novo. Os dois desvios são corrigidos: o estado
estacionário é estável. Duas espécies: cada uma precisa de
$\mu_{i}(S) = D$ para persistir, o que fixa seu próprio $S_{i}^{*}$; o
nutriente só pode assumir um valor, de modo que no máximo uma espécie
persiste. Aquela de menor $S^{*}(D)$ leva $S$ abaixo do que a outra
precisa, e a outra é lavada para fora; qual é essa espécie pode mudar
com $D$, se suas curvas de Monod se cruzarem.
\end{solution}

\begin{solution}{exo:b3:bacteriology:11}
Com a densidade $\rho = N/V$, o estado desligado $A = p_{0}\rho/k$
existe enquanto se mantém abaixo de $A_{\text{thr}}$, isto é,
$\rho < kA_{\text{thr}}/p_{0}$; o estado ligado $A = p_{1}\rho/k$
existe enquanto se mantém acima, isto é, $\rho
> kA_{\text{thr}}/p_{1}$.
Para $kA_{\text{thr}}/p_{1} < \rho <
kA_{\text{thr}}/p_{0}$ os dois são
autoconsistentes: uma população que já ligou se mantém ligada com sua
produção maior, e uma que não ligou fica desligada. Utilidade: o
compromisso --- uma vez que uma colônia tenha decidido construir um
biofilme ou liberar toxina, uma queda de densidade não desfaz a
decisão, e o interruptor é abrupto e sincronizado, em vez de graduado.
\end{solution}

\begin{solution}{exo:b3:bacteriology:12}
(a) Doses altas em tratamentos curtos: o nível fica acima da CIM dos
mutantes e mata igualmente o tipo selvagem e os mutantes de primeiro
passo --- a favor, limitada pela toxicidade. (b) Doses baixas em
tratamentos longos: o nível se instala na janela por dias e cria
resistência --- contra. (c) Parar quando os sintomas cedem: os
sobreviventes são as células menos sensíveis, e o nível em queda passa
pela janela --- contra. (d) Combinação: uma célula precisa ser
resistente aos dois, coisa que um bilhão de células não contém --- a
favor.
\end{solution}

\begin{solution}{pb:b3:bacteriology:1}
\textbf{1.} $18$ gerações: $2^{18} = 2.6\times 10^{5}$ células,
$2.6\times 10^{3}$ por mL.
\textbf{2.} $2\times 10^{9}\times 100 = 2\times 10^{11}$ células $=
2^{37.5}$: depois de \qty{12.5}{h}; massa \qty{0.2}{g}.
\textbf{3.} $6\times 10^{27}/10^{-12} = 6\times 10^{39} = 2^{132}$
células: $132$ gerações, \qty{44}{h}. Os nutrientes, o oxigênio e os
resíduos o detêm em menos de um dia; o crescimento exponencial é sempre
breve.
\textbf{4.} As células estacionárias desmontaram seus ribossomos e
reprimiram suas enzimas metabólicas; precisam reconstruí-los antes de
se dividir. As células exponenciais chegam plenamente equipadas.
\textbf{5.} $\mu = \ln 2/\qty{0.333}{h} = \qty{2.08}{h^{-1}}$; com
$T =
\qty{60}{min}$, \qty{0.69}{h^{-1}}; depois de \qty{6}{h},
$2^{6} = 64$ em vez de $2^{18}$: um fator $2^{12} \approx 4000$ a
menos.
\textbf{6.} $200/(0.1\times 10^{-6}) = 2\times 10^{9}$ células viáveis
por mL. O microscópio também conta células mortas e células que não
crescerão na placa, e um aglomerado dá uma única colônia.
\textbf{7.} $D = 0.5/1 = \qty{0.5}{h^{-1}}$; $T = \ln 2/0.5 =
\qty{1.39}{h} = \qty{83}{min}$.
\textbf{8.} $S^{*} = 0.01\times 0.5/0.5 = \qty{0.01}{g/L}$; $X^{*} =
0.5\times 1.99 \approx \qty{1}{g/L} = 10^{12}$ células por litro,
$10^{9}$ por mL.
\textbf{9.} $0.5\times 10^{12} = 5\times 10^{11}$ células por hora;
$720/1.39
\approx 520$ gerações por mês.
\textbf{10.} $D = 0.95$: $S^{*} = 0.01\times 0.95/0.05 = \qty{0.19}{g/L}$, $X^{*} = 0.5\times 1.81 = \qty{0.9}{g/L}$. Em
$F = \qty{1.1}{L/h}$, $D
> \mu_{\max}$: lavagem.
\textbf{11.} O mutante mantém $S^{*} = 0.01\times 0.5/0.7 = \qty{0.0071}{g/L}$;
nessa glicose a linhagem original cresce a $1.0\times 0.0071/0.0171 =
\qty{0.42}{h^{-1}} < D$ e declina como $e^{-0.08t}$: é lavada para fora
ao longo de algumas semanas. O mutante toma conta.
\textbf{12.} $10^{-9}\times 10^{12} = 1000$ mutações benéficas por
geração: a adaptação é contínua, com um clone vantajoso novo varrendo a
população a cada poucas centenas de gerações --- o quimiostato é uma
máquina de evolução.
\textbf{13.} $kA_{\text{thr}}/p = 0.05\times 6\times 10^{12}/500 =
6\times 10^{8}$ por litro $= 6\times 10^{5}$ células por mL.
\textbf{14.} $10^{10}/6\times 10^{5} \approx 1.7\times 10^{4}$ vezes.
\textbf{15.} $20\times 6\times 10^{12}/500 = 2.4\times 10^{11}$ por
litro, $2.4\times 10^{8}$ por mL --- seis ordens de grandeza acima das
$10^{2}$ por mL do mar: escuro.
\textbf{16.} Dez vezes são $3.3$ duplicações, \qty{6.6}{h}. Com
$10^{9}$ por mL a população ainda está $1700$ vezes acima do limiar: a
luz continua acesa.
\textbf{17.} No mar ninguém vê a luz e ela não alimenta hospedeiro
algum: um mutante escuro poupa \qty{20}{\%} e supera os que brilham. No
órgão, a lula seleciona a favor da luz --- ela expele ou deixa de
alimentar os trapaceiros escuros ---, de modo que o custo compra uma
casa.
\textbf{18.} Bloqueie o receptor ou destrua o \omterm{def:b3:bacteriology:quorum}{autoindutor} (a extinção do
quórum): as bactérias vivem, mas nunca acionam suas toxinas, e o
sistema imune as elimina. Como bactérias sensíveis e ``resistentes''
crescem igualmente --- nada é morto ---, a vantagem seletiva de escapar
do fármaco é pequena, e a resistência deveria se espalhar mais devagar
que a um fármaco que mata.
\textbf{19.} A: $10^{-8}\times 10^{9} = 10$; B: $100$; aos dois:
$10^{-15}\times
10^{9} = 10^{-6}$.
\textbf{20.} $P_{0}(\text{A}) = e^{-10} = 5\times 10^{-5}$;
$P_{0}(\text{ambos}) = e^{-10^{-6}} = 0.999999$.
\textbf{21.} $16 \to 8$: uma meia-vida, \qty{4}{h}; $16 \to 1$: quatro,
\qty{16}{h}; na janela das \qty{4}{h} às \qty{16}{h}, \qty{12}{h} por
dose.
\textbf{22.} Com doses a cada \qty{12}{h}: na janela da hora $4$ à hora
$12$, dois terços do tempo; ao longo de dez dias os mutantes de
primeiro passo de A --- dez deles no início --- são selecionados e A
sozinho falha.
\textbf{23.} Acima de \qty{8}{mg/L} o tempo todo: dose a cada meia-vida
(\qty{4}{h}), ou uma dose com pico de \qty{64}{mg/L} a cada
\qty{12}{h} ($64 \to 8$ em três meias-vidas), ou uma infusão contínua.
Custos: a toxicidade dos picos altos, ou seis doses por dia que os
pacientes não vão cumprir.
\textbf{24.} $10^{-4}\times 10^{9} = 10^{5}$ persistentes sobrevivem
sejam quais forem os fármacos; quando o tratamento para, elas despertam
e regeneram uma população plenamente sensível --- uma recaída. O
tratamento precisa durar o bastante para apanhá-las quando saem da
dormência, e é por isso que a tuberculose leva seis meses.
\textbf{25.} $10^{9}$ células por mL no quimiostato; a luz se acende em
$6\times 10^{5}$ células por mL; $P_{0} = 0.999999$ de que o tratamento
com dois fármacos não encontre célula resistente alguma.
\end{solution}
